
\chapter{Compiler Fundamentals – IRs and Passes}
\section{The Compilation Pipeline}
\begin{itemize}
    \item Front end (lexing, parsing, AST), middle end (IR optimization), back end (code generation).
    \item LLVM IR and GCC GIMPLE: a high‑level view.
\end{itemize}

\begin{figure}[htbp]
\centering
\resizebox{\textwidth}{!}{%
\begin{tikzpicture}[
    node distance=1.7cm,
    box/.style={draw, rounded corners, minimum width=2.8cm, minimum height=0.85cm, align=center},
    smallbox/.style={draw, rounded corners, minimum width=2.4cm, minimum height=0.75cm, align=center},
    arr/.style={-{Latex[length=2mm]}, thick}
]
\node[box] (src) {C/C++ source\\\texttt{atomic\_tx} block};
\node[box, right=of src] (ir) {LLVM IR\\plain loads/stores};
\node[box, right=of ir] (pass) {TM pass\\address policy};
\node[box, right=of pass] (iir) {Instrumented IR\\runtime calls};
\node[box, right=of iir] (obj) {Object code};

\node[smallbox, below=1.2cm of pass] (shared) {\texttt{.tm\_shared}\\registered globals};
\node[smallbox, below=1.2cm of iir] (abi) {TM runtime ABI\\\texttt{tm\_hooks}};

\draw[arr] (src) -- (ir);
\draw[arr] (ir) -- (pass);
\draw[arr] (pass) -- (iir);
\draw[arr] (iir) -- (obj);
\draw[arr] (shared) -- (pass);
\draw[arr] (iir) -- (abi);
\draw[arr] (abi) -- (obj);
\end{tikzpicture}%
}
\caption{A compiler TM pass sits between IR optimization and code generation; its output is ordinary code with calls into the runtime ABI.}
\label{fig:compiler-tm-pass}
\end{figure}

\section{Instrumentation as a Compiler Pass}
\begin{itemize}
    \item Replacing transactional memory accesses with calls to a runtime.
    \item The structure of a \texttt{FunctionPass} in LLVM (and its GIMPLE equivalent).
    \item Pseudocode for the instrumentation algorithm, plus a PlusCal specification.
    \item What a production pass must handle: function calls into uninstrumented code, indirect calls, and escaping pointers.
    \item The instrumentation decision is a static-address-space question: heap/global addresses go through TM; stack addresses bypass it (the \texttt{LLVM\_TM\_ADDR\_CHECK} pattern).
    \item Where TM globals get registered so the runtime can recognise them (the \texttt{.tm\_shared} section of the companion implementation).
    \item Trade-off: instrumenting every load/store (simple, slow) vs.\ annotating shared data (fast, needs programmer discipline).
\end{itemize}

\begin{table}[htbp]
\centering
\begin{tabular}{p{0.22\linewidth}p{0.28\linewidth}p{0.22\linewidth}p{0.16\linewidth}}
\hline
Policy & What the pass rewrites & Programmer\slash{}runtime obligation & Typical cost \\
\hline
Instrument every load/store & All memory operations inside a transaction become runtime calls. & Runtime must cheaply reject private addresses or tolerate over-instrumentation. & Highest; simple and conservative. \\
Static address-space check & Heap and global addresses are instrumented; provable stack addresses bypass TM. & Alias analysis must be conservative when provenance is unclear. & Lower; still safe when uncertain cases are instrumented. \\
Annotated shared section & Only data registered in \texttt{.tm\_shared} is treated as transactional shared state. & Programmer or front end must mark all shared transactional objects. & Lowest common overhead; weakest if annotations are wrong. \\
Manual runtime calls & Source code calls \texttt{tm\_read\_*} and \texttt{tm\_write\_*} directly. & Programmer maintains ABI discipline by hand. & Useful for tests and microbenchmarks; brittle at scale. \\
\hline
\end{tabular}
\caption{Instrumentation policies differ mainly in where the conservative decision is made: compiler, runtime, or programmer.}
\label{tab:instrumentation-policies}
\end{table}

\section{Worked Example: A Load, Before and After the Pass}
\index{instrumentation}\index{LLVM IR}
\label{sec:ir-example}

IR is the right place to see what instrumentation actually changes. A single
transactional read of an \texttt{int} becomes a call into the runtime:

\begin{lstlisting}[language=C, caption={LLVM IR for a transactional load, before and after instrumentation.}]
; before the pass -- a plain memory load
%1 = load i32, i32* %p

; after the pass -- the load is replaced by a runtime call
%1 = call i32 @__tm_read_i4(i32* %p)
\end{lstlisting}

\begin{itemize}
    \item The pass rewrites the memory operand of each load/store into a typed runtime call (\texttt{tm\_read\_i4}, \texttt{tm\_write\_i4}, \ldots).
    \item The static address-space policy decides \emph{which} accesses get rewritten: heap/global go through TM, stack addresses bypass via \texttt{LLVM\_TM\_ADDR\_CHECK}.
    \item The type in the call name matters: \texttt{i4} vs \texttt{i8} select different runtime fast paths and different read-set log widths.
    \item Exercise~1 of this chapter asks you to repeat this mapping by hand on a GIMPLE listing.
    \item The contract the pass instruments \emph{against} --- the exact signatures of \texttt{tm\_read\_*}\slash{}\texttt{tm\_write\_*}, how a transaction begins/ends, and what an abort unwinds --- is the runtime ABI of Chapter~7. If the ABI moves, the pass must move with it; the two are coupled by design, which is why they share the \texttt{tm\_hooks} table rather than a hard-coded call.
\end{itemize}

\section{Worked Example: Instrumenting a Bank Transfer}
\index{bank transfer}\index{runtime ABI}
\label{sec:bank-transfer-instrumented-c}

This listing is not an implementation of NOrec, TL2, or TSC-TM. It is the
shape of the code produced by the pass: ordinary source-level loads and stores
have become calls through the transactional ABI. The stub runtime below makes
the file compile and lets the invariant be tested in a single-threaded run.

\begin{lstlisting}[language=C, caption={Compilable C++17 model of an instrumented bank transfer.}]
#include <assert.h>
#include <stdint.h>

struct Account {
    int balance;
};

extern "C" void __tm_begin(void) {
    /* Real backends allocate per-transaction metadata here. */
}

extern "C" void __tm_end(void) {
    /* Real backends validate and commit here. */
}

extern "C" int __tm_read_i4(const int *p) {
    return *p;
}

extern "C" void __tm_write_i4(int *p, int value) {
    *p = value;
}

static int sum2(const Account *a, const Account *b) {
    return a->balance + b->balance;
}

/* This is the post-pass shape of:
 *
 *   atomic_tx {
 *       if (from->balance - amount >= floor) {
 *           from->balance -= amount;
 *           to->balance += amount;
 *       }
 *   }
 */
static void transfer_instrumented(Account *from,
                                  Account *to,
                                  int amount,
                                  int floor) {
    __tm_begin();

    int from_balance = __tm_read_i4(&from->balance);
    int to_balance   = __tm_read_i4(&to->balance);

    if (from_balance - amount >= floor) {
        __tm_write_i4(&from->balance, from_balance - amount);
        __tm_write_i4(&to->balance,   to_balance + amount);
    }

    __tm_end();
}

int main(void) {
    Account a = { 100 };
    Account b = {  40 };

    int before = sum2(&a, &b);
    transfer_instrumented(&a, &b, 30, 50);
    int after = sum2(&a, &b);

    assert(a.balance >= 50);
    assert(before == after);

    before = after;
    transfer_instrumented(&a, &b, 80, 50);
    after = sum2(&a, &b);

    assert(a.balance >= 50);
    assert(before == after);
    return 0;
}
\end{lstlisting}

Lessons:
\begin{itemize}
    \item The compiler pass does not prove the money-conservation invariant; it preserves the memory operations on which the runtime will enforce atomicity.
    \item The two writes must be in one transaction. Instrumenting each store independently would preserve type correctness but not the bank invariant under concurrency.
    \item The generated calls are typed. Here \texttt{balance} is an \texttt{int}, so the pass emits \texttt{\_\_tm\_read\_i4} and \texttt{\_\_tm\_write\_i4}.
    \item Stack temporaries such as \texttt{from\_balance}, \texttt{to\_balance}, and \texttt{before} are not shared transactional locations and should not be routed through the TM runtime.
\end{itemize}

\section{Worked Example: Model Checking the Address Policy}
\index{TLA+}\index{address-space check}
\label{sec:address-policy-tla}

The pass has a small but important safety rule: do not instrument provably
private stack locations; do instrument heap and global locations. The following
TLA+ module models one pass over a three-instruction program and checks that
the emitted operation stream respects that rule.

\begin{lstlisting}[caption={TLA+ model of the stack-bypass/shared-instrument policy.}]
---- MODULE InstrumentPass ----
EXTENDS TLC, Sequences, Naturals

VARIABLES pc, out

Ops == {"load", "store"}
Addrs == {"stack", "heap", "global"}
Kinds == {"plain", "tmread", "tmwrite"}

Prog ==
  << [op |-> "load",  addr |-> "stack"],
     [op |-> "store", addr |-> "heap"],
     [op |-> "load",  addr |-> "global"] >>

Instr(op, addr) ==
  IF addr = "stack" THEN
      "plain"
  ELSE IF op = "load" THEN
      "tmread"
  ELSE
      "tmwrite"

Init ==
  /\ pc = 1
  /\ out = << >>

Next ==
  IF pc <= Len(Prog) THEN
      /\ out' = Append(out, Instr(Prog[pc].op, Prog[pc].addr))
      /\ pc' = pc + 1
  ELSE
      /\ UNCHANGED << pc, out >>

TypeOK ==
  /\ pc \in 1..(Len(Prog) + 1)
  /\ out \in Seq(Kinds)
  /\ Len(out) <= Len(Prog)

StackBypass ==
  \A i \in 1..Len(out):
    Prog[i].addr = "stack" => out[i] = "plain"

SharedInstrumented ==
  \A i \in 1..Len(out):
    Prog[i].addr # "stack" =>
      ((Prog[i].op = "load"  => out[i] = "tmread") /\
       (Prog[i].op = "store" => out[i] = "tmwrite"))

Spec == Init /\ [][Next]_<< pc, out >>

====
\end{lstlisting}

Lessons:
\begin{itemize}
    \item The model is deliberately smaller than LLVM. It isolates the address-policy obligation from parsing, dominance, and code generation.
    \item \texttt{StackBypass} catches accidental over-instrumentation of private stack slots.
    \item \texttt{SharedInstrumented} catches the dangerous bug: a heap or global access escaping the TM runtime.
    \item In a production pass, the hard case is not the rule itself but proving that an address is stack-private after optimization and aliasing.
\end{itemize}

\section{Summary}
\begin{itemize}
    \item Instrumentation is a source‑to‑source (or IR‑to‑IR) transformation.
    \item A pass that instruments every memory access is the simplest possible TM compiler support.
    \item A production-quality pass needs an address-space policy, a shared-data registry, and careful handling of uninstrumented callees.
\end{itemize}

\section{Exercises}
\begin{enumerate}
    \item Given a small C function, manually identify which GIMPLE statements would be replaced by calls to a TM runtime barrier.
    \item Critique a naive instrumentation pass that instruments all loads and stores in the whole program. What performance problems would arise?
    \item Design an LLVM pass that instruments only loads and stores inside basic blocks marked with a special intrinsic.
\end{enumerate}
